\hypertarget{v11-coupled_structure_full-constant-linear-coordinates-and-coupled-dilation}{%
\section{Constant linear coordinates and coupled
dilation}\label{v11-coupled_structure_full-constant-linear-coordinates-and-coupled-dilation}}

\hypertarget{v11-coupled_structure_full-exact-meaning-of-decoupling}{%
\subsection{Exact meaning of
decoupling}\label{v11-coupled_structure_full-exact-meaning-of-decoupling}}

Let \(q,r,a,b>0\), with contraction \(<1\) imposed when discussing the
normalized entire solution. For an arbitrary analytic vector input
\(X=(u,v)\) with \(u(0)=v(0)=0\), define the functional vector field

\[
\mathcal F[X](z)=
\begin{pmatrix}1+u(qz)v(rz)\\1+u(az)v(bz)\end{pmatrix}.
\tag{S1}
\]

A nonzero constant row \(c=(c_1,c_2)\in\mathbb C^2\) is an
\textbf{autonomous scalar projection} if some scalar functional
\(\mathcal H\) satisfies

\[
c\mathcal F[X]=\mathcal H[cX]
\quad\text{for every such analytic input }X.
\tag{S2}
\]

The functional \(\mathcal H\) may use the whole scalar input. Thus this
definition allows any scalar dilation equation, not only a prescribed
quadratic ansatz. Equality (S2) means that two vector inputs with the
same scalar projection have the same projected vector field. It is an
identity of functional differential operators. It does not assert an
identity only along the distinguished normalized solution.

A constant invertible change \(W=PX\) completely decouples the operator
if both rows of \(P\) are autonomous scalar projections. No assumption
about the form of their scalar functionals is needed.

\hypertarget{v11-coupled_structure_full-classification-of-autonomous-linear-projections}{%
\subsection{Classification of autonomous linear
projections}\label{v11-coupled_structure_full-classification-of-autonomous-linear-projections}}

\textbf{Theorem S1.} For positive nonzero scales, a nonzero row \(c\)
satisfies (S2) if and only if

\[
q=a,\qquad r=b,\qquad c=\gamma(1,-1),\qquad\gamma\ne0.
\tag{S3}
\]

In that exceptional case the projected vector field is identically zero.
Consequently, no invertible constant linear change completely decouples
(S1), including when some or all scales coincide. If \((q,r)\ne(a,b)\),
there is not even one nonzero autonomous scalar projection. Four
pairwise distinct contraction scales are a special case of this stronger
result.

\textbf{Proof.} Subtract the constant forcing. Define

\[
Q_c[u,v](z)=c_1u(qz)v(rz)+c_2u(az)v(bz).
\]

If (S2) holds, \(Q_c\) is constant on every fiber of
\((u,v)\mapsto c_1u+c_2v\). Fix \(z_0\ne0\). Polynomial interpolation
gives a scalar polynomial \(f\) with \(f(0)=0\) and \(f(sz_0)=1\) at
every distinct scale \(s\in\{q,r,a,b\}\). For inputs \(u=xf,v=yf\),
where \(x,y\in\mathbb C\),

\[
Q_c[xf,yf](z_0)=(c_1+c_2)xy.
\tag{S4}
\]

It must depend only on \(c_1x+c_2y\). If \(c_1+c_2\ne0\),
differentiating the polynomial \(xy\) along the fiber direction
\((c_2,-c_1)\) would give

\[
c_2y-c_1x=0\quad\text{for every }x,y,
\]

which forces \(c=0\), a contradiction. Therefore \(c_2=-c_1\), and
scaling the row reduces the question to \(c=(1,-1)\).

For this row, put \(u=w+g,v=g\). This input and \((w,0)\) have the same
projection \(w\). The second input has \(Q_c[w,0]=0\), so autonomy
forces

\[
w(qz)g(rz)-w(az)g(bz)
+g(qz)g(rz)-g(az)g(bz)=0
\tag{S5}
\]

for all zero-constant analytic \(w,g\). Setting \(w=0\) and subtracting
that instance of (S5) leaves

\[
w(qz)g(rz)=w(az)g(bz).
\tag{S6}
\]

Taking successively \((w,g)=(z,z),(z^2,z),(z,z^2)\) yields \(qr=ab\),
\(q^2r=a^2b\), and \(qr^2=ab^2\). Nonzero scales then give \(q=a\) and
\(r=b\). Conversely these equalities make the two components of (S1)
identical on every input, so \((1,-1)\mathcal F[X]=0\). The only
possible autonomous rows therefore span a one-dimensional space, which
cannot be the two rows of an invertible matrix. \(\square\)

The same proof works on odd analytic inputs for positive scales: choose
\(f(z)=z p(z^2)\) by interpolation at the distinct positive scale
squares, and use \((z,z),(z^3,z),(z,z^3)\) in (S6). Positivity then
resolves the equalities of squares. The theorem does not depend on
allowing even inputs.

\hypertarget{v11-coupled_structure_full-tensor-explanation-for-four-distinct-scales}{%
\subsection{Tensor explanation for four distinct
scales}\label{v11-coupled_structure_full-tensor-explanation-for-four-distinct-scales}}

The preceding proof also makes a coordinate invariant obstruction
explicit. Under an invertible state change, write the original
coordinates as \(u=\ell_1(W),v=\ell_2(W)\), where \(\ell_1,\ell_2\) are
independent linear forms. At the pair of distinct nodes \((qz,rz)\), the
first quadratic term has bilinear coefficient

\[
B=\ell_1\otimes\ell_2.
\]

It has rank one, but is not symmetric:
\(B-B^{\mathsf T}=\ell_1\wedge\ell_2\ne0\). At four distinct nonzero
nodes, polynomial interpolation lets their vector values vary
independently, so terms supported at \((az,bz)\) cannot cancel this
coefficient. A scalar autonomous channel at the same node pair would
require \(B=\kappa\ell\otimes\ell\), a symmetric rank-one tensor, after
projecting the output onto that channel. For any output row with a
nonzero coefficient on this term, the alternating part remains nonzero.
The other node pair gives the same obstruction if the first output
coefficient vanishes.

This is not a rank-only argument: the tensor \(\ell_1\otimes\ell_2\)
already has rank one. Its nonzero alternating part is the decisive
invariant at distinct nodes. When the nodes of a product coincide, the
symmetric quadratic form \(uv\) has matrix

\[
\begin{pmatrix}0&1/2\\1/2&0\end{pmatrix},
\]

of rank two, whereas a square of one linear coordinate has rank one.
Rank is then the appropriate obstruction. The fiber proof in Theorem S1
handles node coincidences and cancellations without mixing these two
cases.

\hypertarget{v11-coupled_structure_full-particular-solutions-and-degenerate-parameters}{%
\subsection{Particular solutions and degenerate
parameters}\label{v11-coupled_structure_full-particular-solutions-and-degenerate-parameters}}

Theorem S1 is an operator statement. It does not exclude coincident
components for particular initial data, algebraic relations along a
solution, elimination into a scalar equation of higher order, nonlinear
coordinate changes, or identities with other named special functions.

For contraction scales, the normalized analytic germ is unique. If
\((a,b)=(q,r)\), its components coincide because their derivatives are
identical and their initial values agree. The common component is the
scalar family \(T'=1+T(qz)T(rz),T(0)=0\). The same coincidence holds if
\((a,b)=(r,q)\): \((T,T)\) solves both equations because multiplication
is commutative, and uniqueness identifies the normalized solution. For
\(q\ne r\), this reversed-pair case has no autonomous scalar operator
projection by Theorem S1, despite having \(U=V\) along its normalized
solution. This is a concrete reason the two notions must remain
distinct.

In fact, equality of normalized components occurs exactly when the
unordered scale pairs agree:

\[
U=V\quad\Longleftrightarrow\quad\{q,r\}=\{a,b\}.
\tag{S7}
\]

To prove necessity, write \(U=V=T=z+c_1z^3+\cdots\). Then
\(3c_1=qr=ab>0\). Comparing the \(z^4\) coefficients of
\(T(qz)T(rz)=T(az)T(bz)\) gives \(c_1qr(q^2+r^2)=c_1ab(a^2+b^2)\), so
both products and sums of squares agree. Positive scales then have equal
sums and are the roots of the same quadratic. Sufficiency is the
uniqueness argument above. For example, \((q,r,a,b)=(1/2,1/3,1/3,1/2)\)
has coincident normalized components and no nonzero autonomous operator
projection.

At \(q=r=a=b=s>0\), the variables \(w=u+v,d=u-v\) give

\[
w'=2+\tfrac12\bigl(w(sz)^2-d(sz)^2\bigr),\qquad d'=0.
\]

This triangular system is not a completely decoupled operator. The
normalized solution lies in the invariant subspace \(d=0\).

Zero scales lie outside Theorem S1. On zero-constant inputs, a product
vanishes identically if either of its scales is zero. If both products
vanish, both equations are already \(u'=v'=1\). If only the first
product vanishes, the projection \(u\) is autonomous, regardless of
whether the second pair agrees with the first. Thus extending Theorem S1
to zero scales would be false. Scale one and positive scales greater
than one do not affect the algebraic operator theorem, although the
contraction-domain entire growth theorems do not automatically extend to
them. Negative or complex scales and odd-input coincidences caused by
signs are not classified here.

\hypertarget{v11-coupled_structure_full-claim-status}{%
\subsection{Claim status}\label{v11-coupled_structure_full-claim-status}}

Theorem S1 and (S7) are proved above. The proof is elementary and
independent of the coupled growth theorem. Historical priority for these
precise operator statements is not asserted. The theorem provides a
specified nondecoupling result, not a universal assertion that the
normalized functions cannot reduce to any previously studied function
class.
